\cleardoublepage
\thispagestyle{empty}
\begin{center}


{\Large\textbf{Photon – Theorie und Anwendungen}}\\[1.2em]
{\large Dipl.-Ing.\,(FH)\,Christian Weilharter}\\[1.2em]
\textcopyright~2025, Christian Weilharter, Traunstein\\[2em]
\end{center}
\begin{flushleft}
	\begin{tabular}{@{}l l}
	
	
		\textbf{ISBN: (Print)} & 978-3-912302-09-7 \\[0.5em]
			\textbf{ISBN: (E-Book)} & 978-3-912302-01-1 \\[0.5em]
		\textbf{Auflage:} & 1.~Auflage 2025 \\[0.5em]
		\textbf{Satz:} & \LaTeX \\[0.5em]
		\textbf{Verlag:} & Christian Weilharter \\[0.5em]
		\textbf{Druck:} & Amazon KDP (Print-on-Demand) \\[0.5em]

			\textbf{E-Book-Ausgabe:} & Apple Books\\[0.5em]
			&Amazon Kindle Direct Publishing \\[0.5em]
		\textbf{Kontakt:} & \href{mailto:info@mathandphysics.de}{info@mathandphysics.de}\\[0.5em]
		\textbf{Web:} & \href{https://www.mathandphysics.de}{www.mathandphysics.de}\\
		

	\end{tabular}
\end{flushleft}

\vspace{2em}
\noindent
Alle Rechte vorbehalten. Kein Teil dieses Buches darf ohne schriftliche Genehmigung des Autors 
in irgendeiner Form reproduziert, gespeichert oder übertragen werden, 
weder elektronisch, mechanisch, durch Fotokopien, Aufnahmen noch auf andere Weise.
\begin{center}\small Printed in Germany\end{center}

\cleardoublepage

\chapter*{Vorwort}
\markboth{Vorwort}{Vorwort} % sorgt für die Kopfzeile, optional
% kein addcontentsline → Vorwort erscheint NICHT im Inhaltsverzeichnis



Die Entstehung dieses Buches war von einer tiefen Faszination für das Licht und seinem fundamentalen Vermittler – dem Photon – getragen. In der modernen Physik nimmt das Photon eine zentrale Rolle ein: als Teilchen ohne Masse, aber mit Energie und Impuls; als Bote der elektromagnetischen Wechselwirkung; und als Schlüsselfigur der Quantenmechanik.
	\noindent
Mein Ziel war es, dieses vielseitige Konzept in seiner historischen Entwicklung, seiner physikalischen Bedeutung und seinen technischen Anwendungen so darzustellen, dass sowohl interessierte Laien als auch fortgeschrittene Leser einen fundierten Zugang erhalten – ohne unnötige Vereinfachung, aber stets verständlich.
	\noindent
Bei der Konzeption und Ausarbeitung dieses Buches kam auch moderne Technologie zum Einsatz. Zur Formulierung, Strukturierung und Reflexion von Inhalten wurde das KI-System ChatGPT von OpenAI (Stand 2025) unterstützend eingesetzt. Diese Form der Zusammenarbeit ermöglichte es, Gedanken schneller zu konkretisieren, alternative Formulierungen zu prüfen und komplexe Zusammenhänge noch klarer darzustellen.
Alle Inhalte wurden jedoch von mir kritisch geprüft, überarbeitet und verantwortet.
	\noindent
Dieses Buch ist somit nicht nur ein Beitrag zur Didaktik der modernen Physik – es ist auch ein Experiment, wie traditionelle Wissenschaft mit modernen Werkzeugen eine neue Form der Klarheit und Zugänglichkeit gewinnen kann.
	\noindent
Ich hoffe, dass die Begeisterung für das Thema auf die Leserinnen und Leser überspringt – so wie sie mich seit Jahren begleitet.



\begin{flushright}
	\textit{Dipl.-Ing.(FH) Christian Weilharter} \\
	\vspace{0.5em}
	[Traunstein, 2025]
\end{flushright}


\chapter*{Aufbau des Buches}
\markboth{Aufbau des Buches}{Aufbau des Buches} % sorgt für die Kopfzeile, optional
% kein addcontentsline → Vorwort erscheint NICHT im Inhaltsverzeichnis

Dieses Buch ist in acht thematisch aufeinander abgestimmte Kapitel gegliedert. Es führt den Leser von den physikalischen Grundlagen des Lichts über die Quantisierung des elektromagnetischen Feldes bis hin zu modernen Anwendungen und offenen Fragen der Forschung. Jedes Kapitel verfolgt das Ziel, die Rolle des Photons aus einer bestimmten Perspektive zu beleuchten – historisch, theoretisch, experimentell oder technologisch.

\begin{itemize}
	\item \textbf{Kapitel I} beschreibt kurz den Überblick über die historische Entwicklung  und Theorie des Photons.
	\item \textbf{Kapitel II} schildert den Weg von der klassischen Lichttheorie zur Einführung des Lichtquants. Dabei werden zentrale Experimente wie der Photoeffekt\index{Photoeffekt} und die Schwarzkörperstrahlung\index{Schwarzkörperstrahlung} behandelt, die zur Quantentheorie führten.
	
	\item \textbf{Kapitel III} widmet sich den physikalischen Eigenschaften des Photons, darunter Energie, Impuls, Spin, Polarisation und der Aspekt seiner Masselosigkeit.\index{Polarisation}
	
	\item \textbf{Kapitel IV} beleuchtet die experimentelle Bestätigung des Photons – von Millikans Photoeffektmessung\index{Millikan, Robert A.} bis hin zu modernen Einzelphotonenexperimenten und Quantenradierern\index{Quantenradierer}.
	
	\item \textbf{Kapitel V} führt in die Quantenelektrodynamik (QED) ein und zeigt, wie das Photon dort als Eichboson der elektromagnetischen Wechselwirkung verstanden wird.
	
	\item \textbf{Kapitel VI} beschreibt die praktischen Anwendungen des Photons, etwa in der Lasertechnologie, der medizinischen Bildgebung und der Kommunikationstechnik.\index{Medizinische Bildgebung}
	
	\item \textbf{Kapitel VII} zeigt aktuelle Entwicklungen und Zukunftspotenziale auf, insbesondere im Bereich der Quanteninformation,\index{Quanteninformation} Photonik und fundamentalen Forschung.
	
	\item \textbf{Kapitel VIII} ordnet das Photon in das Standardmodell der Teilchenphysik ein.\index{Standardmodell der Teilchenphysik} Es wird gezeigt, wie das masselose Photon aus Symmetrien der Theorie hervorgeht und welche offenen Fragen sich daraus ergeben.
\end{itemize}
Zahlreiche Abbildungen, Zitate und Reflexionen begleiten die fachlichen Inhalte, um auch philosophische, historische und erkenntnistheoretische Dimensionen sichtbar zu machen. Die mathematischen Mittel werden behutsam eingeführt, sodass das Werk auch für interessierte Leser mit naturwissenschaftlichem Grundverständnis zugänglich bleibt.

\begin{quote}
	\textit{Mit dem Verständnis für die historische, konzeptionelle und experimentelle Basis des Photons im Gepäck, betreten wir im nächsten Kapitel das Terrain der Quantisierung – jenen Schritt, der das Licht aus der klassischen Welt in die Quantenphysik überführt.}
\end{quote}

\subsubsection*{Hinweise zum Lesen}
Um die unterschiedlichen Aspekte der Darstellung klar voneinander zu trennen,
werden im gesamten Buch farbige Boxen eingesetzt. Sie dienen der Orientierung
und helfen dabei, zentrale Inhalte hervorzuheben:

\begin{itemize}
	\item \textbf{Physikboxen} (blau) geben physikalische Erklärungen und
	Hintergrundinformationen.
	\item \textbf{Matheboxen} (grün) enthalten Herleitungen, Formeln und
	mathematische Details.
	\item \textbf{Didaktikboxen} (gelb) weisen auf Denkfallen hin oder liefern
	alternative Erklärungen.
	\item \textbf{Hinweisboxen} (grau) geben Strukturhinweise oder Verweise auf
	andere Kapitel und Anhänge.
	\item \textbf{Warnboxen} (rot) markieren besonders kritische Aspekte oder
	Missverständnisse.
	\item \textbf{Hypoboxen} (orange) beleuchten hypothetische Szenarien oder
	„Was-wäre-wenn“-Fragen.
\end{itemize}


% -------------------------------------------------------------
\section*{Hinweis zur Open-Access-Version}
% -------------------------------------------------------------

Dieses Buch ist Teil der Open-Science-Initiative „Christian \& Co-Pilot – Math \& Physics“.

Die vollständige, farbige PDF-Version ist frei zugänglich unter:

\begin{itemize}
	\item \href{https://mathandphysics.eu}{\texttt{https://mathandphysics.de}}
	\item \href{https://zenodo.org/communities/christian-copilot}{\texttt{https://zenodo.org/communities/christian-copilot}}
\end{itemize}
Die gedruckte Ausgabe wurde für komfortables Lesen, langfristige Archivierung
und den Aufbau einer Open-Access-Bibliothek in Printform erstellt.
Mit dem Kauf dieses Buches unterstützen Sie die freie wissenschaftliche Publikation
und den Gedanken einer offenen, überprüfbaren Wissenschaft.